\subsection{Residual covers and strict continuum screening}
A cellwise residual bracket need not solve another action problem at every state. On each rectangular state cell, convex interpolation of the vertex upper values dominates $V_1$; each affine function \eqref{eq:affine95} minorizes it globally. The same interpolation and subgradient construction applies to convex $f$. Comparing an affine minorant of one function with the convex interpolation upper bound of the other at every vertex proves a bound for $V_1-f$ throughout the cell. Taking the best available anchor and then the worst cell gives a whole-domain residual interval. The $32\times32$ cover below therefore qualifies the continuous square, rather than converting its $1,089$ nodes into empirical generalization observations.

Let $[l_i,u_i]$ enclose an original successor value over every point of bin $B_i$. If $\tau$ is the $k$th largest $l_i$, $k=n(1-\alpha)$, then
\begin{equation}\label{eq:strict95}
 u_i<\tau
\end{equation}
excludes all risk mass from bin $i$ in an upper-tail maximizer. At least $1-\alpha$ of the original probability mass has value at least $\tau$; strict separation makes any positive mass on the excluded bin suboptimal. Equality is retained. These bounds must be available before screening. The screening calculation below pays for their construction and for the test itself, and compares identical upper-tail values with and without screening. This continuum-bin result is separate from enumerating all extreme kernels of a finite risk envelope.

